\section{Introduction}

A tracker that has lost a person and finds them again must decide whether
the return belongs to the old identity. ByteTrack
\cite{zhang2022bytetrack} keeps a lost track alive for thirty frames and
re-attaches any later detection that clears its association threshold.
When the decision is right, an identity survives an occlusion and the
trajectory continues. When it is wrong, nothing announces the error: the
wrong person inherits the identifier, every later frame is attached to
that identifier, and every count, dwell time or trajectory built on top of
it is wrong from that moment on. Reconnection failures are silent, which
is why they deserve more attention than their share of association events
would suggest.

Work on the problem takes two forms. One improves the association cost
itself, through motion models, matching cascades and observation-centric
updates \cite{zhang2022bytetrack,aharon2022botsort,cao2023ocsort}. The
other adds appearance, either as a re-identification branch inside the
tracker or as a linking stage after it. We ask a narrower question that
neither answers directly. Given the proposals a tracker has already made,
which ones should be committed?

Three deployment settings make that narrow question worth asking. A
privacy-constrained installation may not permit storing face or clothing
embeddings at all. An edge device may not have the budget to run a second
network beside the detector. A retrospective analysis may inherit a tracker
whose association rule cannot be changed, only wrapped. In each case the
available lever is not a better match but a better decision about whether to
commit the match, and in each case the decision has to be made from the
evidence at that frame.

We treat reactivation as a decision and place a scalar gate in front of
it. The gate reads the detector confidence, the predicted box residuals,
the track age, the observed gap, and the margins between the proposed
assignment and its best competitors in the same row and column. A
regularized logistic score combines these, and product terms let the
weight of motion agreement and of assignment ambiguity change with gap
length. A refused detection returns to the tracker as unmatched, so the
decision changes the state that later frames see. No appearance network is
trained and no gallery is stored: the gate decides from geometry.

The paper makes four claims. \textbf{First, the effect survives matched
admission counts.} Counting failures rewards silence, so we re-select
thresholds until the native association cost, an additive score and the
interaction score admit the same number of actions per source. At the cap
of twelve admissions per source, ByteTrack then fails in 98 of 381
evaluated reactivations against 158 of 390 for its native cost, and
verified-correct reactivations rise from 232 to 283. \textbf{Second, the
comparison is against a published score, not only a hand-set threshold.}
After porting the BoostTrack association similarity into the same host and
checking the port element by element, the interaction score fails less
often in all five operating points we tested, by 9.72 to 13.50 percentage
points. \textbf{Third, the gain comes from the product terms and not from
extra summaries of the competing assignments}: a score that drops the
products and keeps only the three summaries has no contrast interval that
excludes zero at any of the five operating points, while the full score
lowers the failure proportion by 3.29 points at the point where the effect
is established. \textbf{Fourth, the cost is
explicit.} Relative to an unconstrained ByteTrack the gate removes 1,397
evaluated failures and 136 natural gaps, and HOTA falls by 1.65 points. We
report that as the price of the decision rather than argue it away.

One boundary applies throughout. The gate uses no appearance signal, so
the claim covers detection-only pipelines; and every failure proportion in
this paper is an observation on the split that names it, not a bound on
unseen sources.
